\documentclass[aspectratio=169]{beamer}
\usepackage[utf8]{inputenc}
\usepackage[french]{babel}
\usepackage{graphicx}
\usepackage{listings}
\usepackage{xcolor}
\usepackage{tikz}
\usepackage{booktabs}
\usepackage{fontawesome5}
\usepackage{hyperref}

% Theme
\usetheme{Madrid}
\usecolortheme{default}

% Couleurs personnalisées
\definecolor{bleuNgoa}{RGB}{0,102,204}
\definecolor{vertSucces}{RGB}{0,153,51}
\definecolor{rougeErreur}{RGB}{204,0,0}
\definecolor{grisCode}{RGB}{248,248,248}
\definecolor{orangeWarning}{RGB}{255,140,0}

\setbeamercolor{structure}{fg=bleuNgoa}
\setbeamercolor{frametitle}{bg=bleuNgoa,fg=white}

% Configuration des listings (code)
\lstset{
    backgroundcolor=\color{grisCode},
    basicstyle=\ttfamily\footnotesize,
    keywordstyle=\color{blue},
    commentstyle=\color{gray},
    stringstyle=\color{red},
    showstringspaces=false,
    breaklines=true,
    frame=single,
    numbers=left,
    numberstyle=\tiny\color{gray}
}

% Informations du document
\title[SmartSearch]{SmartSearch}
\subtitle{Système Intelligent de Gestion de Documents Perdus}
\author{NGOA (HERVEMEUPS)}
\institute{M2 SIGL Professionnel}
\date{Juin 2026}

\begin{document}

% ============================================================
% SLIDE 1 : PAGE DE TITRE
% ============================================================
\begin{frame}
    \titlepage
    \begin{center}
        \Large
        \textbf{Analyse, Conception et Implémentation} \\
        \large
        d'un Système Intelligent Dédié à la Déclaration \\
        et à la Recherche de Documents Perdus

        \vspace{0.5cm}
        \normalsize
        \textbf{Version 3.1} - Juin 2026
    \end{center}
\end{frame}

% ============================================================
% SLIDE 2 : TABLE DES MATIÈRES
% ============================================================
\begin{frame}{Plan de la Présentation}
    \tableofcontents
\end{frame}

% ============================================================
% SECTION 1 : CONTEXTE & PROBLÉMATIQUE
% ============================================================
\section{Contexte \& Problématique}

\begin{frame}{Contexte : La Perte de Documents}
    \begin{columns}[T]
        \begin{column}{0.48\textwidth}
            \begin{block}{Constat}
                \begin{itemize}
                    \item Perte fréquente de documents officiels
                    \item CNI, Passeports, Permis, Diplômes...
                    \item Démarches administratives longues
                    \item Coûts de remplacement élevés
                \end{itemize}
            \end{block}

            \vspace{0.3cm}

            \begin{alertblock}{Problème}
                \textbf{Aucun système centralisé} pour signaler et rechercher des documents perdus
            \end{alertblock}
        \end{column}
        \begin{column}{0.48\textwidth}
            \begin{block}{Besoins Identifiés}
                \begin{enumerate}
                    \item Déclarer un document perdu
                    \item Déclarer un document trouvé
                    \item Rechercher des correspondances
                    \item Notifier automatiquement
                    \item Gérer les utilisateurs
                \end{enumerate}
            \end{block}

            \vspace{0.3cm}

            \begin{exampleblock}{Solution}
                Plateforme web intelligente avec matching automatique
            \end{exampleblock}
        \end{column}
    \end{columns}
\end{frame}

\begin{frame}{Problématique}
    \begin{center}
        \Large
        \textbf{Comment créer un système intelligent capable de détecter automatiquement les correspondances entre documents perdus et trouvés, même avec des informations partielles ou erronées ?}
    \end{center}

    \vspace{1cm}

    \begin{block}{Défis Techniques}
        \begin{itemize}
            \item Matching avec informations incomplètes
            \item Tolérance aux fautes de frappe
            \item Performance en temps réel
            \item Scalabilité
            \item Sécurité des données personnelles
        \end{itemize}
    \end{block}
\end{frame}

% ============================================================
% SECTION 2 : ARCHITECTURE GLOBALE
% ============================================================
\section{Architecture Globale}

\begin{frame}{Architecture Technique}
    \begin{center}
        \begin{tikzpicture}[
            box/.style={rectangle, draw, fill=blue!10, text width=2.5cm, align=center, minimum height=0.8cm},
            arrow/.style={->, >=stealth, thick}
        ]
            % Frontend
            \node[box, fill=green!20] (frontend) at (0,4) {\textbf{Frontend} \\ HTML/CSS/JS};

            % Backend
            \node[box, fill=blue!20] (backend) at (5,4) {\textbf{Backend} \\ Node.js + Express};

            % AI Service
            \node[box, fill=purple!20] (ai) at (10,4) {\textbf{Service IA} \\ Python + FastAPI};

            % Base de données
            \node[box, fill=orange!20] (db) at (5,2) {\textbf{Base de Données} \\ MongoDB Atlas};

            % Cloud
            \node[box, fill=red!20] (cloud) at (5,0.5) {\textbf{Hébergement} \\ Render.com};

            % Flèches
            \draw[arrow] (frontend) -- node[above] {\tiny API REST} (backend);
            \draw[arrow] (backend) -- node[right] {\tiny Matching} (ai);
            \draw[arrow] (backend) -- node[left] {\tiny CRUD} (db);
            \draw[arrow, dashed] (frontend) -- (db);
            \draw[arrow] (db) -- (cloud);
            \draw[arrow] (backend) -- (cloud);
            \draw[arrow] (ai) -- (cloud);
        \end{tikzpicture}
    \end{center}

    \vspace{0.3cm}

    \begin{block}{Architecture 3-tiers}
        \textbf{Présentation} (Frontend) → \textbf{Logique} (Backend + IA) → \textbf{Données} (MongoDB)
    \end{block}
\end{frame}

\begin{frame}{Stack Technologique}
    \begin{columns}[T]
        \begin{column}{0.32\textwidth}
            \begin{block}{Frontend}
                \begin{itemize}
                    \item HTML5 / CSS3
                    \item JavaScript Vanilla
                    \item Font Awesome
                    \item Chart.js
                \end{itemize}
            \end{block}
        \end{column}
        \begin{column}{0.32\textwidth}
            \begin{block}{Backend}
                \begin{itemize}
                    \item Node.js 18+
                    \item Express.js
                    \item JWT (Auth)
                    \item Bcrypt
                    \item Fuse.js
                \end{itemize}
            \end{block}
        \end{column}
        \begin{column}{0.32\textwidth}
            \begin{block}{IA \& Data}
                \begin{itemize}
                    \item Python 3.9+
                    \item FastAPI
                    \item RapidFuzz
                    \item Transformers
                    \item MongoDB
                \end{itemize}
            \end{block}
        \end{column}
    \end{columns}

    \vspace{0.5cm}

    \begin{exampleblock}{Choix Technologiques}
        \begin{itemize}
            \item \textbf{JavaScript} : Écosystème riche, performance
            \item \textbf{Python} : IA/ML, bibliothèques NLP
            \item \textbf{MongoDB} : NoSQL, flexibilité, scalabilité
            \item \textbf{Render} : Gratuit, CI/CD automatique
        \end{itemize}
    \end{exampleblock}
\end{frame}

% ============================================================
% SECTION 3 : FONCTIONNALITÉS
% ============================================================
\section{Fonctionnalités}

\begin{frame}{Fonctionnalités Utilisateur}
    \begin{columns}[T]
        \begin{column}{0.48\textwidth}
            \begin{block}{1. Authentification}
                \begin{itemize}
                    \item Inscription / Connexion
                    \item JWT (24h)
                    \item Rôles : Admin / Déclarant
                    \item Sécurité : Bcrypt
                \end{itemize}
            \end{block}

            \vspace{0.3cm}

            \begin{block}{2. Déclaration}
                \begin{itemize}
                    \item Type : Perte / Découverte
                    \item Documents : CNI, Passeport, Permis...
                    \item Informations partielles acceptées
                    \item Localisation (ville, quartier)
                \end{itemize}
            \end{block}
        \end{column}
        \begin{column}{0.48\textwidth}
            \begin{block}{3. Recherche Intelligente}
                \begin{itemize}
                    \item Fuzzy matching
                    \item Recherche par critères
                    \item Scoring de pertinence
                    \item Tri par score
                \end{itemize}
            \end{block}

            \vspace{0.3cm}

            \begin{block}{4. Matching Automatique}
                \begin{itemize}
                    \item Détection auto PERTE ↔ DECOUVERTE
                    \item Notification utilisateurs
                    \item Confiance : Faible → Très Haute
                    \item Dashboard des correspondances
                \end{itemize}
            \end{block}
        \end{column}
    \end{columns}
\end{frame}

\begin{frame}{Fonctionnalités Admin}
    \begin{block}{Dashboard Administrateur}
        \begin{columns}[T]
            \begin{column}{0.48\textwidth}
                \textbf{Gestion Utilisateurs}
                \begin{itemize}
                    \item CRUD complet
                    \item Réinitialisation password
                    \item Statistiques utilisateurs
                    \item Bannir / Réactiver
                \end{itemize}
            \end{column}
            \begin{column}{0.48\textwidth}
                \textbf{Gestion Déclarations}
                \begin{itemize}
                    \item Validation
                    \item Suppression
                    \item Statistiques
                    \item Export de données
                \end{itemize}
            \end{column}
        \end{columns}
    \end{block}

    \vspace{0.3cm}

    \begin{block}{Analytics}
        \begin{itemize}
            \item Graphiques interactifs (Chart.js)
            \item Évolution temporelle
            \item Documents par type / lieu
            \item Taux de récupération
        \end{itemize}
    \end{block}
\end{frame}

% ============================================================
% SECTION 4 : SYSTÈME DE MATCHING
% ============================================================
\section{Système de Matching Intelligent}

\begin{frame}{Moteur de Matching Multi-Niveaux}
    \begin{center}
        \begin{tikzpicture}[
            box/.style={rectangle, draw, fill=blue!10, text width=2.8cm, align=center, minimum height=1cm},
            arrow/.style={->, >=stealth, thick}
        ]
            % Titre
            \node[box, fill=bleuNgoa!30, text width=11cm] (titre) at (0,3.5) {
                \textbf{Moteur de Matching Intelligent}
            };

            % Niveau 1
            \node[box, fill=green!20] (fuzzy) at (-4,2) {
                \textbf{Niveau 1} \\ Fuzzy Matching \\ (RapidFuzz) \\ \small 40\%
            };

            % Niveau 2
            \node[box, fill=blue!20] (embed) at (0,2) {
                \textbf{Niveau 2} \\ Embeddings \\ (Transformers) \\ \small 40\%
            };

            % Niveau 3
            \node[box, fill=purple!20] (llm) at (4,2) {
                \textbf{Niveau 3} \\ LLM \\ (Claude 3) \\ \small 20\%
            };

            % Scoring
            \node[box, fill=orange!20, text width=11cm] (score) at (0,0.5) {
                \textbf{Score Global} = 0.4×Fuzzy + 0.4×Embeddings + 0.2×LLM
            };

            % Décision
            \node[box, fill=red!20, text width=11cm] (decision) at (0,-0.8) {
                \textbf{Décision} : Score ≥ 0.65 → \textcolor{vertSucces}{MATCH} | Score < 0.65 → \textcolor{rougeErreur}{NO MATCH}
            };

            % Flèches
            \draw[arrow] (fuzzy) -- (score);
            \draw[arrow] (embed) -- (score);
            \draw[arrow] (llm) -- (score);
            \draw[arrow] (score) -- (decision);
        \end{tikzpicture}
    \end{center}
\end{frame}

\begin{frame}{Niveau 1 : Fuzzy Matching avec RapidFuzz}
    \begin{block}{Algorithmes Utilisés}
        \begin{enumerate}
            \item \texttt{fuzz.ratio()} - Similarité globale (Levenshtein)
            \begin{itemize}
                \item "KOUAKOU" vs "KOAKOU" → 92.86\%
            \end{itemize}

            \item \texttt{fuzz.partial\_ratio()} - Détection de sous-chaînes
            \begin{itemize}
                \item "123456" in "CI0123456789" → 100\%
            \end{itemize}

            \item \texttt{fuzz.token\_sort\_ratio()} - Ordre des mots
            \begin{itemize}
                \item "Jean KOUASSI" vs "KOUASSI Jean" → 100\%
            \end{itemize}
        \end{enumerate}
    \end{block}

    \begin{exampleblock}{Bonus Progressif}
        ≥90\% → +0.25 | ≥80\% → +0.20 | ≥70\% → +0.15 | ≥60\% → +0.10
    \end{exampleblock}
\end{frame}

\begin{frame}{Niveau 2 : Embeddings Sémantiques}
    \begin{block}{Sentence Transformers}
        \begin{itemize}
            \item Modèle : \texttt{paraphrase-multilingual-MiniLM-L12-v2}
            \item Conversion texte → vecteur (384 dimensions)
            \item Similarité cosinus entre embeddings
            \item Compréhension du sens, pas juste des mots
        \end{itemize}
    \end{block}

    \vspace{0.3cm}

    \begin{exampleblock}{Exemple}
        \begin{itemize}
            \item "CNI perdue à Abidjan" → [0.12, -0.45, 0.78, ...]
            \item "Carte identité trouvée Abidjan" → [0.15, -0.42, 0.81, ...]
            \item Similarité cosinus → 0.87 (très similaire)
        \end{itemize}
    \end{exampleblock}
\end{frame}

\begin{frame}{Niveau 3 : LLM (Claude 3 Sonnet)}
    \begin{block}{Analyse Contextuelle Avancée}
        \begin{itemize}
            \item Compréhension du contexte global
            \item Raisonnement sur les correspondances
            \item Validation des matchs fuzzy/embeddings
            \item Génération d'explications
        \end{itemize}
    \end{block}

    \vspace{0.3cm}

    \begin{exampleblock}{Exemple de Raisonnement LLM}
        \small
        \textit{"Forte correspondance probable : même type de document (CNI), nom très similaire (KOUAKOU/KOAKOU - probablement une faute de frappe), numéro partiel correspond, même localisation (Abidjan), dates compatibles (écart de 2 jours). Confiance : HAUTE."}
    \end{exampleblock}
\end{frame}

% ============================================================
% SECTION 5 : SÉCURITÉ
% ============================================================
\section{Sécurité \& Protection des Données}

\begin{frame}{Mesures de Sécurité}
    \begin{columns}[T]
        \begin{column}{0.48\textwidth}
            \begin{block}{Authentification}
                \begin{itemize}
                    \item JWT (JSON Web Token)
                    \item Expiration 24h
                    \item Refresh token
                    \item Hachage Bcrypt (10 rounds)
                \end{itemize}
            \end{block}

            \vspace{0.3cm}

            \begin{block}{Autorisation}
                \begin{itemize}
                    \item Gestion des rôles (Admin/User)
                    \item Middleware d'authentification
                    \item Validation des permissions
                    \item CORS configuré
                \end{itemize}
            \end{block}
        \end{column}
        \begin{column}{0.48\textwidth}
            \begin{block}{Protection des Données}
                \begin{itemize}
                    \item Variables d'environnement (.env)
                    \item Secrets non versionnés
                    \item HTTPS obligatoire
                    \item Validation des entrées
                \end{itemize}
            \end{block}

            \vspace{0.3cm}

            \begin{block}{Conformité}
                \begin{itemize}
                    \item Minimisation des données
                    \item Données partielles (nom/numéro)
                    \item Pas de stockage complet de documents
                    \item Suppression possible
                \end{itemize}
            \end{block}
        \end{column}
    \end{columns}
\end{frame}

% ============================================================
% SECTION 6 : BASE DE DONNÉES
% ============================================================
\section{Base de Données}

\begin{frame}[fragile]{Modèle de Données MongoDB}
    \begin{columns}[T]
        \begin{column}{0.48\textwidth}
            \begin{block}{Collection Users}
                \small
\begin{lstlisting}[language=json]
{
  "_id": ObjectId,
  "username": "string",
  "email": "string",
  "password": "hash",
  "role": "ADMIN|DECLARANT",
  "dateInscription": Date
}
\end{lstlisting}
            \end{block}
        \end{column}
        \begin{column}{0.48\textwidth}
            \begin{block}{Collection Declarations}
                \small
\begin{lstlisting}[language=json]
{
  "_id": ObjectId,
  "type": "PERTE|DECOUVERTE",
  "typeDocument": "CNI|...",
  "nomPartiel": "string",
  "numeroPartiel": "string",
  "description": "string",
  "localisation": {
    "ville": "string",
    "quartier": "string"
  },
  "declarantId": ObjectId,
  "date": Date
}
\end{lstlisting}
            \end{block}
        \end{column}
    \end{columns}

    \vspace{0.3cm}

    \begin{exampleblock}{Avantages NoSQL}
        Schéma flexible, indexation efficace, scalabilité horizontale
    \end{exampleblock}
\end{frame}

% ============================================================
% SECTION 7 : RÉSULTATS & PERFORMANCE
% ============================================================
\section{Résultats \& Performance}

\begin{frame}{Performance du Système}
    \begin{table}
        \begin{tabular}{lcc}
            \toprule
            \textbf{Métrique} & \textbf{Valeur} & \textbf{Cible} \\
            \midrule
            Temps de réponse API & <200ms & <500ms \\
            Temps de matching & <500ms & <1000ms \\
            Throughput & 200+ req/sec & 100 req/sec \\
            Disponibilité & 99.5\% & 99\% \\
            Taux de réussite tests & 100\% (15/15) & 95\% \\
            \bottomrule
        \end{tabular}
    \end{table}

    \vspace{0.5cm}

    \begin{exampleblock}{Performance Fuzzy Matching}
        \begin{itemize}
            \item \textbf{RapidFuzz} : >10,000 comparaisons/seconde
            \item \textbf{Embeddings} : Cache intelligent (Redis)
            \item \textbf{LLM} : Appels asynchrones, early stopping
        \end{itemize}
    \end{exampleblock}
\end{frame}

\begin{frame}{Impact du Fuzzy Matching}
    \begin{center}
        \begin{tikzpicture}
            % Avant
            \draw[fill=rougeErreur!30] (0,0) rectangle (3,4);
            \node at (1.5,4.5) {\textbf{AVANT (v3.0)}};

            \node[text width=2.5cm, align=center] at (1.5,3) {
                \small Correspondances détectées
            };
            \node at (1.5,2) {\Huge 65\%};
            \node at (1.5,0.5) {\Large \textcolor{rougeErreur}{\faTimesCircle}};

            % Flèche
            \draw[->, ultra thick] (3.5,2) -- (5,2);
            \node[above, text width=1.2cm, align=center] at (4.25,2) {\tiny RapidFuzz v3.1};

            % Après
            \draw[fill=vertSucces!30] (5.5,0) rectangle (8.5,4);
            \node at (7,4.5) {\textbf{APRÈS (v3.1)}};

            \node[text width=2.5cm, align=center] at (7,3) {
                \small Correspondances détectées
            };
            \node at (7,2) {\Huge 95\%};
            \node at (7,0.5) {\Large \textcolor{vertSucces}{\faCheckCircle}};

            % Gain
            \node[fill=orangeWarning, text width=8cm, align=center] at (4.25,-0.5) {
                \textbf{Gain : +30 points} → +35\% de correspondances
            };
        \end{tikzpicture}
    \end{center}
\end{frame}

% ============================================================
% SECTION 8 : DÉPLOIEMENT
% ============================================================
\section{Déploiement \& Production}

\begin{frame}{Infrastructure de Production}
    \begin{block}{Hébergement Render.com}
        \begin{columns}[T]
            \begin{column}{0.48\textwidth}
                \textbf{Backend Node.js}
                \begin{itemize}
                    \item Plan Free
                    \item Auto-deploy GitHub
                    \item Variables d'environnement
                    \item Logs en temps réel
                \end{itemize}
            \end{column}
            \begin{column}{0.48\textwidth}
                \textbf{Service IA Python}
                \begin{itemize}
                    \item Plan Free
                    \item Build automatique
                    \item requirements.txt
                    \item Cold start : ~30s
                \end{itemize}
            \end{column}
        \end{columns}
    \end{block}

    \vspace{0.3cm}

    \begin{block}{MongoDB Atlas}
        \begin{itemize}
            \item Cluster M0 (gratuit)
            \item 512 MB storage
            \item Réplication automatique
            \item Backup journalier
        \end{itemize}
    \end{block}
\end{frame}

\begin{frame}[fragile]{Pipeline CI/CD}
    \begin{center}
        \begin{tikzpicture}[
            box/.style={rectangle, draw, fill=blue!10, text width=2cm, align=center, minimum height=0.8cm},
            arrow/.style={->, >=stealth, thick}
        ]
            % GitHub
            \node[box, fill=black!10] (git) at (0,0) {\faGithub \\ GitHub};

            % Render
            \node[box, fill=blue!20] (render) at (3,0) {Render \\ (Webhook)};

            % Build
            \node[box, fill=orange!20] (build) at (6,0) {Build \\ \& Test};

            % Deploy
            \node[box, fill=green!20] (deploy) at (9,0) {Deploy \\ Production};

            % Flèches
            \draw[arrow] (git) -- node[above] {\tiny push} (render);
            \draw[arrow] (render) -- node[above] {\tiny trigger} (build);
            \draw[arrow] (build) -- node[above] {\tiny success} (deploy);
        \end{tikzpicture}
    \end{center}

    \vspace{0.5cm}

    \begin{block}{Processus Automatique}
        \small
\begin{lstlisting}[language=bash]
1. git push origin main
2. Render webhook detecte push
3. Build: npm install / pip install
4. Tests unitaires (si presents)
5. Deploy automatique
6. Health check endpoint
\end{lstlisting}
    \end{block}
\end{frame}

% ============================================================
% SECTION 9 : TESTS & VALIDATION
% ============================================================
\section{Tests \& Validation}

\begin{frame}{Stratégie de Tests}
    \begin{columns}[T]
        \begin{column}{0.48\textwidth}
            \begin{block}{Tests Unitaires}
                \begin{itemize}
                    \item Fuzzy matching (15 tests)
                    \item Contrôleurs backend
                    \item Services métier
                    \item Utilitaires
                \end{itemize}
                \textbf{Taux de réussite : 100\%}
            \end{block}

            \vspace{0.3cm}

            \begin{block}{Tests d'Intégration}
                \begin{itemize}
                    \item API endpoints
                    \item Authentification JWT
                    \item CRUD déclarations
                    \item Matching automatique
                \end{itemize}
            \end{block}
        \end{column}
        \begin{column}{0.48\textwidth}
            \begin{block}{Tests Manuels}
                \begin{itemize}
                    \item Scénarios utilisateur
                    \item Interface graphique
                    \item Compatibilité navigateurs
                    \item Responsive design
                \end{itemize}
            \end{block}

            \vspace{0.3cm}

            \begin{block}{Tests de Performance}
                \begin{itemize}
                    \item Charge (wrk)
                    \item Stress testing
                    \item Temps de réponse
                    \item Scalabilité
                \end{itemize}
            \end{block}
        \end{column}
    \end{columns}
\end{frame}

\begin{frame}[fragile]{Exemple de Test Fuzzy Matching}
    \begin{block}{Test Python}
        \small
\begin{lstlisting}[language=Python]
def test_fuzzy_matching_with_typo():
    """Test detection de faute de frappe"""
    result = fuzz.ratio("KOUAKOU", "KOAKOU")
    assert result >= 90, f"Expected >= 90%, got {result}%"

def test_partial_number_match():
    """Test numero partiel"""
    result = fuzz.partial_ratio("123456", "CI0123456789")
    assert result == 100, f"Expected 100%, got {result}%"

def test_city_with_typo():
    """Test ville avec faute"""
    result = fuzz.ratio("Abidjan", "Abijan")
    assert result >= 85, f"Expected >= 85%, got {result}%"
\end{lstlisting}
    \end{block}

    \begin{exampleblock}{Résultat}
        \textcolor{vertSucces}{\faCheckCircle} \textbf{15/15 tests réussis (100\%)}
    \end{exampleblock}
\end{frame}

% ============================================================
% SECTION 10 : DÉMONSTRATION
% ============================================================
\section{Démonstration}

\begin{frame}{Scénario Utilisateur Complet}
    \begin{enumerate}
        \item \textbf{Inscription} : Créer un compte (Alice)
        \begin{itemize}
            \item POST /api/auth/register
        \end{itemize}

        \vspace{0.2cm}

        \item \textbf{Déclaration de Perte} : Alice perd sa CNI
        \begin{itemize}
            \item POST /api/declarations
            \item Type: PERTE, Nom: "KOUAKOU", Numéro: "123456"
        \end{itemize}

        \vspace{0.2cm}

        \item \textbf{Inscription} : Créer un compte (Bob)

        \vspace{0.2cm}

        \item \textbf{Déclaration de Découverte} : Bob trouve une CNI (avec fautes)
        \begin{itemize}
            \item POST /api/declarations
            \item Type: DECOUVERTE, Nom: "KOAKOU", Numéro: "CI0123456789"
        \end{itemize}

        \vspace{0.2cm}

        \item \textbf{Matching Automatique} : Le système détecte la correspondance
        \begin{itemize}
            \item Score: 0.86 → \textcolor{vertSucces}{\faCheckCircle} MATCH
            \item Notification Alice \& Bob
        \end{itemize}
    \end{enumerate}
\end{frame}

\begin{frame}{Interface Utilisateur}
    \begin{block}{Pages Principales}
        \begin{enumerate}
            \item \textbf{Accueil} : Présentation, statistiques
            \item \textbf{Connexion / Inscription}
            \item \textbf{Déclaration} : Formulaire guidé
            \item \textbf{Recherche} : Filtres multiples, résultats scorés
            \item \textbf{Mes Déclarations} : Historique personnel
            \item \textbf{Dashboard Admin} : Gestion complète
        \end{enumerate}
    \end{block}

    \vspace{0.3cm}

    \begin{exampleblock}{Design}
        \begin{itemize}
            \item Responsive (mobile-first)
            \item Moderne et épuré
            \item Accessibilité (WCAG 2.1)
            \item Thème cohérent
        \end{itemize}
    \end{exampleblock}
\end{frame}

% ============================================================
% SECTION 11 : DIFFICULTÉS & SOLUTIONS
% ============================================================
\section{Difficultés Rencontrées}

\begin{frame}{Défis Techniques}
    \begin{block}{1. Matching avec Informations Incomplètes}
        \textbf{Problème} : Utilisateurs ne connaissent que des fragments \\
        \textbf{Solution} : Fuzzy matching + embeddings sémantiques
    \end{block}

    \begin{block}{2. Performance du LLM}
        \textbf{Problème} : Appels LLM lents (>2s) \\
        \textbf{Solution} : Cache, early stopping, fallback NLP
    \end{block}

    \begin{block}{3. Déploiement Gratuit}
        \textbf{Problème} : Limitations plans gratuits (cold start) \\
        \textbf{Solution} : Render + MongoDB Atlas, optimisation build
    \end{block}

    \begin{block}{4. Sécurité des Données}
        \textbf{Problème} : Données sensibles (documents d'identité) \\
        \textbf{Solution} : Stockage partiel uniquement, JWT, HTTPS
    \end{block}
\end{frame}

% ============================================================
% SECTION 12 : PERSPECTIVES
% ============================================================
\section{Perspectives d'Évolution}

\begin{frame}{Améliorations Futures}
    \begin{columns}[T]
        \begin{column}{0.48\textwidth}
            \begin{block}{Court Terme (3 mois)}
                \begin{itemize}
                    \item Notifications email/SMS
                    \item Application mobile
                    \item Export PDF déclarations
                    \item Multi-langues (EN, FR)
                \end{itemize}
            \end{block}

            \vspace{0.3cm}

            \begin{block}{Moyen Terme (6 mois)}
                \begin{itemize}
                    \item Machine Learning (patterns)
                    \item Matching phonétique
                    \item Géolocalisation GPS
                    \item Système de réputation
                \end{itemize}
            \end{block}
        \end{column}
        \begin{column}{0.48\textwidth}
            \begin{block}{Long Terme (12 mois)}
                \begin{itemize}
                    \item OCR (extraction documents)
                    \item IA prédictive (lieux)
                    \item Blockchain (traçabilité)
                    \item API publique
                \end{itemize}
            \end{block}

            \vspace{0.3cm}

            \begin{block}{Partenariats}
                \begin{itemize}
                    \item Commissariats de police
                    \item Mairies / Préfectures
                    \item Compagnies de transport
                    \item Établissements publics
                \end{itemize}
            \end{block}
        \end{column}
    \end{columns}
\end{frame}

% ============================================================
% SECTION 13 : CONCLUSION
% ============================================================
\section{Conclusion}

\begin{frame}{Bilan du Projet}
    \begin{block}{Objectifs Atteints}
        \begin{itemize}
            \item \textcolor{vertSucces}{\faCheckCircle} Plateforme web fonctionnelle et déployée
            \item \textcolor{vertSucces}{\faCheckCircle} Système de matching intelligent (+95\% de détection)
            \item \textcolor{vertSucces}{\faCheckCircle} Interface utilisateur intuitive
            \item \textcolor{vertSucces}{\faCheckCircle} Dashboard administrateur complet
            \item \textcolor{vertSucces}{\faCheckCircle} Sécurité et protection des données
            \item \textcolor{vertSucces}{\faCheckCircle} Documentation exhaustive
            \item \textcolor{vertSucces}{\faCheckCircle} Tests validés (100\%)
        \end{itemize}
    \end{block}

    \vspace{0.3cm}

    \begin{exampleblock}{Compétences Mobilisées}
        Full-stack (JS, Python), IA/ML (NLP, LLM), DevOps (CI/CD), UX/UI, Gestion de projet
    \end{exampleblock}
\end{frame}

\begin{frame}{Apports du Projet}
    \begin{columns}[T]
        \begin{column}{0.48\textwidth}
            \begin{block}{Techniques}
                \begin{itemize}
                    \item Architecture 3-tiers
                    \item Microservices (Backend + IA)
                    \item Fuzzy matching avancé
                    \item Intégration LLM
                    \item NoSQL (MongoDB)
                    \item CI/CD automatisé
                \end{itemize}
            \end{block}
        \end{column}
        \begin{column}{0.48\textwidth}
            \begin{block}{Personnels}
                \begin{itemize}
                    \item Gestion de projet
                    \item Problem solving
                    \item Veille technologique
                    \item Documentation
                    \item Autonomie
                    \item Rigueur
                \end{itemize}
            \end{block}
        \end{column}
    \end{columns}

    \vspace{0.5cm}

    \begin{alertblock}{Impact Social}
        Faciliter la récupération de documents perdus pour les citoyens
    \end{alertblock}
\end{frame}

\begin{frame}{Conclusion Générale}
    \begin{center}
        \Large
        \textbf{SmartSearch} est une plateforme web intelligente qui transforme la recherche de documents perdus grâce à un système de \textbf{matching automatique} basé sur l'\textbf{intelligence artificielle}.

        \vspace{1cm}

        Le système atteint \textbf{95\% de détection} des correspondances, même avec des \textbf{informations partielles ou erronées}, grâce à une combinaison innovante de \textbf{fuzzy matching}, \textbf{embeddings sémantiques} et \textbf{LLM}.

        \vspace{1cm}

        \huge
        \textbf{Merci pour votre attention !}
    \end{center}
\end{frame}

% ============================================================
% QUESTIONS
% ============================================================
\begin{frame}
    \begin{center}
        \Huge
        \textbf{Questions ?}

        \vspace{2cm}

        \normalsize
        \textbf{NGOA (HERVEMEUPS)} \\
        M2 SIGL Professionnel - Juin 2026 \\
        \vspace{0.5cm}
        SmartSearch - Version 3.1 \\
        Système Intelligent de Gestion de Documents Perdus

        \vspace{1cm}

        \small
        \faGithub\ GitHub : HERVEMEUPS/smartsearch-app \\
        \faGlobe\ URL : smartsearch-backend.onrender.com
    \end{center}
\end{frame}

% ============================================================
% ANNEXES
% ============================================================
\appendix
\section{Annexes}

\begin{frame}{Annexe : Statistiques du Projet}
    \begin{table}
        \small
        \begin{tabular}{lc}
            \toprule
            \textbf{Métrique} & \textbf{Valeur} \\
            \midrule
            Lignes de code (Backend) & ~5,000 \\
            Lignes de code (Frontend) & ~3,000 \\
            Lignes de code (IA Service) & ~2,000 \\
            Lignes de documentation & ~10,600 \\
            Nombre de commits & 30+ \\
            Durée du projet & 3 mois \\
            Nombre de fichiers & 80+ \\
            Tests automatisés & 15 \\
            \bottomrule
        \end{tabular}
    \end{table}
\end{frame}

\begin{frame}{Annexe : Technologies \& Versions}
    \begin{table}
        \footnotesize
        \begin{tabular}{lll}
            \toprule
            \textbf{Technologie} & \textbf{Version} & \textbf{Usage} \\
            \midrule
            Node.js & 18+ & Backend API \\
            Express.js & 4.18 & Framework web \\
            MongoDB & 6.0 & Base de données \\
            Python & 3.9+ & Service IA \\
            FastAPI & 0.109 & API IA \\
            RapidFuzz & 3.6.1 & Fuzzy matching \\
            Transformers & 4.36 & Embeddings \\
            Claude API & 3.0 & LLM \\
            JWT & 9.0 & Authentification \\
            Bcrypt & 5.1 & Hachage passwords \\
            \bottomrule
        \end{tabular}
    \end{table}
\end{frame}

\begin{frame}{Annexe : Ressources \& Documentation}
    \begin{block}{Documentation Projet}
        \begin{itemize}
            \item README.md - Vue d'ensemble
            \item ARCHITECTURE.md - Architecture technique
            \item FUZZY\_MATCHING\_SYSTEM.md - Système de matching
            \item GUIDE\_TEST\_FUZZY\_MATCHING.md - Tests
            \item DEPLOYMENT\_GUIDE.md - Déploiement
            \item CHANGELOG.md - Historique des versions
        \end{itemize}
    \end{block}

    \begin{block}{Ressources Externes}
        \begin{itemize}
            \item RapidFuzz : \url{https://github.com/maxbachmann/RapidFuzz}
            \item MongoDB Atlas : \url{https://www.mongodb.com/cloud/atlas}
            \item Render : \url{https://render.com}
            \item Anthropic Claude : \url{https://www.anthropic.com}
        \end{itemize}
    \end{block}
\end{frame}

\end{document}
